\scanpage{448}
\mspage{360}
accessible stable par $\underrightarrow{\lim}$ filtrantes, $\pi$ un cardinal utile pour $M$. Alors pour tt $x\in\mathrm{Ob}\,M$, on a une représentation
\[
x=\underset{i\unsure{\in|\notin} I}{\underrightarrow{\lim}}\,\unsure{x_i|x.}
\]
$L$-$\pi$-adaptée, [i.e. avec 1) $I$ équiv. à une petite catégorie et $I$ grde devant $\pi$, 2) $I$ stable par $\underrightarrow{\lim}$ filtrantes de cardin\unsure{al|aux} $\leq\pi$, 3) $I\unsure{\text{---}|\longrightarrow}M$ commute à ces $\underrightarrow{\lim}$, \unsure{4)|1)|d)} $x_i\in\mathrm{Ob}\,M_\pi\quad\forall i\in I$], avec de plus $I$ \uline{ordonnée}.

En effet, on part de
\[
x=\underset{j\in J}{\underrightarrow{\lim}}\,x_j,\qquad
J\cong(M_\pi/x)_0
\]
\inserted{\unsure{une} sous-catégorie réduite équiv. à $M_{\pi/x}$}

\deleted{\unsure{On pr}} On définit $I$ comme l'ens des sous-\unsure{catéories|catégories} \uline{filtrantes} \deleted{\unsure{simples}} \inserted{$J'$} de $J$, de \deleted{\unsure{cardinal}} \inserted{\deleted{\unsure{card}}} $\leq\pi$, \deleted{Il est imm} $I$ ordonné par l'inclusion \unsure{?|$\supseteq$|$\subseteq$}. \deleted{Il a} et on pose, pour $J'\in I$
\[
x_{J'}=\underset{j\in J'}{\underrightarrow{\lim}}\,x_j.
\]
Il est immédiat qu'on obtient ainsi un système inductif dans $M$, indexé par $I$.
% archive sheet number: 447
